\documentclass[11pt]{article}

\usepackage{amsmath}
\usepackage{amsfonts}
\usepackage{amsthm,thmtools}
\usepackage{tikz}
\usepackage{enumerate}


\newcommand{\batsym}{\begin{tikzpicture}[scale=0.05]
	\draw[fill,black]   (-0.25,1.48) .. controls (-0.1,1.5) and (0.1,1.5) .. (0.25,1.48) -- (0.35,1.92) .. controls (0.425,1.8) and (0.41,1.3) .. (0.45,1.2) .. controls (0.6,1.05) and (1.96,1.05) .. (1.98,2.08) -- (5.93,2.08) .. controls (4.2,1.45) and (4,0.3) .. (4.2,-0.28) .. controls (2.4,-0.09) and (0.4,-0.5) .. (0,-2.052) .. controls (-0.4,-0.5) and (-2.4,-0.09) .. (-4.2,-0.28) .. controls (-4,0.3) and (-4.2,1.45) .. (-5.93,2.08) -- (-1.98,2.08) .. controls (-1.96,1.05) and (-0.6,1.05) .. (-0.45,1.2) .. controls (-0.41,1.3) and (-0.425,1.8) .. (-0.35,1.92) -- (-0.25,1.48);
	\end{tikzpicture}}
\newcommand{\QEB}{\begin{flushright}\batsym\end{flushright}}



\def\bold{\textbf}
\theoremstyle{plain}
\newtheorem{problem}{Problem}
\theoremstyle{remark}
\newtheorem*{idea}{Idea}

\renewcommand*{\proofname}{Solution}


\renewcommand{\qedsymbol}{$\batsym$}

\newcommand {\ds}{\displaystyle\sum}
\newcommand {\nnR} {\mathbb{R}_{\geq 0}}

\newcommand {\Mtt}[4] {\begin{bmatrix}#1 & #2\\ #3 & #4\end{bmatrix}}
\newcommand {\Ctt}[2] {\Mtt {#1}{-#2}{#2}{#1}}
\newcommand {\mbR} {\mathbb R}
\newcommand {\mcC} {\mathcal C}
\newcommand {\mbZ} {\mathbb Z}
\newcommand {\mbN} {\mathbb N}
\newcommand {\mbQ} {\mathbb Q}
\newcommand {\RnQ} {\mathbb R \setminus \mathbb Q}

\begin{document}
	\title{Math 3A03 - Tutorial 3 Questions - Winter 2019}
	\author{Nikolay Hristov}
	\date{January 28/30, 2019}
	\maketitle
	
	\begin{problem}
		Prove, using the algebra of limits, $\displaystyle \lim_{n\rightarrow \infty}\left(\frac 1 {n^3}+\frac{n}{2(n+3)}\right)= \frac 1 2$.
	\end{problem}

	\begin{problem}
		Suppose that $a_n\geq 0$. Prove if $\displaystyle \lim_{n\rightarrow \infty} a_n \rightarrow L$ then $L\geq 0$ and $\displaystyle \lim_{n\rightarrow \infty} \sqrt{a_n} \rightarrow \sqrt{L}$.
	\end{problem}

\begin{problem}
	Let $F_1=1, F_2=1, F_n=F_{n-1}+F_{n-2}$,
	
	\begin{enumerate}[a)]	
		\item Prove $F_n$ diverges as $n\rightarrow \infty$.
		
		\item How can you modify the above to show $F_n$ diverges to $\infty$?
	\end{enumerate}
\end{problem}

\begin{problem}
	Let $x_n\rightarrow \infty$ (i.e. diverges to $\infty$), then $s_n=\frac{x_n}{x_n+1}$ converges to 1 as $n\rightarrow\infty$.
\end{problem}


\begin{problem}
	Show that the definition of the limit of a sequence is equivalent to the following slight modification (note this isn't the one from your assignment!):
	
	A sequence $\left\{s_n\right\}$ converges to L if given any $\epsilon>0$ there is some real number $K$ such that if $n\geq K$ then $\left| s_n-L \right|<\epsilon$.
\end{problem}

\begin{problem}
	
	Show that ``the set $E$ is dense in $\mathbb R$'' is equivalent to ``given any $y\in \mathbb R$ there exists a sequence of points $x_n\in E$ such that $\displaystyle \lim_{n \rightarrow \infty} x_n=y$''.
\end{problem}
	
\end{document}